\section{Normalized Gaussian wave functions}
\label{sec:states}
\begin{definition}[Schwartz space]\label{def:schwartz}
Let $\mathcal D=\Sch$ be the Schwartz space: a smooth function $f$ belongs to
$\mathcal D$ when
\[
 \sup_{x\in\R}|x|^r|f^{(k)}(x)|<\infty
 \qquad (r,k\in\mathbb N_0).
\]
\end{definition}
We use Lebesgue measure and put $\mathcal H=L^2(\R;\C)$. Elements of $\mathcal H$ are equivalence classes of functions: two functions
represent the same element if they agree almost everywhere. Schwartz decay makes
every $f\in\mathcal D$ square-integrable. The usual inclusion
\[
 J:\mathcal D\longrightarrow\mathcal H,\qquad Jf=[f],
\]
sends a Schwartz function to its class. It is injective: continuous functions
that agree almost everywhere agree everywhere. We henceforth identify $f$ with $Jf$
when taking an $L^2$ norm or inner product, writing
\[
 \|f\|_2=\|Jf\|_{\mathcal H},\qquad
 \langle f,g\rangle=\int_\R\overline{f(x)}g(x)\,dx.
\]
The inner product is conjugate-linear in its first argument.

\begin{definition}[Normalized representative]\label{def:state}
A normalized state representative is a Schwartz function $f\in\mathcal D$
with $\|f\|_2=1$.
\end{definition}

For $\hbar>0$, position and momentum act on the common invariant domain
$\mathcal D$ by
\[
 (Xf)(x)=xf(x),\qquad (P_\hbar f)(x)=-i\hbar f'(x).
\]
Multiplication by $x$ and differentiation preserve Schwartz space. Their images
therefore have well-defined $L^2$ norms and inner products.

\begin{definition}[Means and covariance]\label{def:covariance}
For an operator
$A:\mathcal D\to\mathcal D$ and a normalized $f$, define
\[
 \mu_f(A)=\Rea\langle f,Af\rangle,\qquad
 u_f(A)=Af-\mu_f(A)f.
\]
For the ordered pair $(R_1,R_2)=(X,P_\hbar)$, the covariance is
\begin{equation}\label{eq:state-covariance}
 (\gamma_f)_{jk}=2\Rea\langle u_f(R_j),u_f(R_k)\rangle.
\end{equation}
\end{definition}
In particular, $(\gamma_f)_{jj}=2\|u_f(R_j)\|_2^2$ is twice the variance.

\subsection{Existence and normalization}
For a width $s>0$, we consider the positive real function
\begin{equation}\label{eq:gaussian}
 \psi_s(x)=\frac{1}{\sqrt{\sqrt\pi\,s}}
            \exp\left(-\frac{x^2}{2s^2}\right).
\end{equation}
\begin{theorem}\label{thm:normalized}
For every $s>0$, the function $\psi_s$ belongs to $\mathcal D$ and satisfies
$\|\psi_s\|_2=1$. It determines a unique normalized state representative
with the pointwise values in~\eqref{eq:gaussian}.
\end{theorem}
\begin{proof}
Put $\phi(x)=e^{-x^2/2}$. We define the \emph{Hermite polynomials} by
\[
 H_0(x)=1,\qquad H_{k+1}(x)=xH_k(x)-H_k'(x)\quad(k\geq0).
\]
The \emph{Hermite derivative identity} is
\[
 \phi^{(k)}(x)=(-1)^k H_k(x)\phi(x).
\]
It follows by induction from $\phi'=-x\phi$ and the product rule.
For every $r\geq0$, $|x|^r\phi(x)\to0$ as $|x|\to\infty$.
Continuity bounds this function globally. Expanding $H_k$ into monomials
then bounds $|x|^r|\phi^{(k)}(x)|$ for each $r,k$.
Dilation by $x\mapsto x/s$, inclusion of real values into $\C$, and
multiplication by a constant preserve Schwartz space, so $\psi_s\in\mathcal D$.

The Gaussian integral gives
\[
 \|\psi_s\|_2^2
 =\frac{1}{\sqrt\pi\,s}\int_\R e^{-x^2/s^2}\,dx
 =\frac{\sqrt\pi\,s}{\sqrt\pi\,s}=1.
\]
A function with the prescribed pointwise values is equal to $\psi_s$.
\end{proof}
The pointwise formula~\eqref{eq:gaussian} fixes the overall phase of $\psi_s$.

\subsection{Action and moments}
\begin{theorem}\label{thm:moments}
For every $s,\hbar>0$,
\begin{gather}
 P_\hbar\psi_s=\frac{i\hbar}{s^2}X\psi_s,
 \qquad \mu_{\psi_s}(X)=\mu_{\psi_s}(P_\hbar)=0,\label{eq:action}\\
 \|X\psi_s\|_2^2=\frac{s^2}{2},\qquad
 \|P_\hbar\psi_s\|_2^2=\frac{\hbar^2}{2s^2},\label{eq:seconds}\\
 \langle X\psi_s,P_\hbar\psi_s\rangle=\frac{i\hbar}{2}.
 \label{eq:cross}
\end{gather}
The first identity holds as an equality of Schwartz functions.
\end{theorem}
\begin{proof}
Differentiating~\eqref{eq:gaussian} yields $\psi_s'=-(x/s^2)\psi_s$ and hence
the momentum action. The integrand in
$\langle \psi_s,X\psi_s\rangle$ is odd and integrable, so this complex
inner product vanishes.

Set $A=\|X\psi_s\|_2^2$. With $\langle\cdot,\cdot\rangle_\R$
denoting the real part of the complex inner product, integration by parts gives
\begin{equation}\label{eq:ibp}
 1-\frac{A}{s^2}
 =\langle \psi_s,(X\psi_s)'\rangle_\R
 =-\langle \psi_s',X\psi_s\rangle_\R
 =\frac{A}{s^2}.
\end{equation}
The first equality uses $(X\psi_s)'=\psi_s+x\psi_s'$ and normalization.
Schwartz decay removes the boundary term, and the derivative formula gives
the last equality. Thus $A=s^2/2$.

Complex linearity in the second argument now gives
\[
 \langle \psi_s,P_\hbar\psi_s\rangle=0,\qquad
 \langle X\psi_s,P_\hbar\psi_s\rangle
 =\frac{i\hbar}{s^2}A=\frac{i\hbar}{2}.
\]
Finally, $\|P_\hbar\psi_s\|_2^2=(\hbar^2/s^4)A=\hbar^2/(2s^2)$.
\end{proof}

\subsection{Covariance and units}
\begin{theorem}\label{thm:covariance}
For every $s,\hbar>0$, the covariance of $\psi_s$ is
\[
 \gamma_{\psi_s}=\begin{pmatrix}s^2&0\\0&\hbar^2/s^2\end{pmatrix}.
\]
If $T=\diag(s^{-1},s/\hbar)$ and
$\widetilde\gamma=T\gamma_{\psi_s}T^T$, then $\widetilde\gamma=I_2$
and $\widetilde\gamma$ is admissible in the sense
of~\eqref{eq:admissibility}.
\end{theorem}
\begin{proof}
The means vanish, so the centered vectors in~\eqref{eq:state-covariance}
are $X\psi_s$ and $P_\hbar\psi_s$. Equations~\eqref{eq:seconds}
give the diagonal entries. The cross moment~\eqref{eq:cross} is purely
imaginary, as is its complex conjugate, so both off-diagonal entries vanish.
Diagonal multiplication gives $\widetilde\gamma=I_2$.
Proposition~\ref{prop:diagonal}, with $a=b=1$, proves admissibility.
\end{proof}

The matrix $T$ converts position and momentum to the dimensionless quantities
$X/s$ and $sP_\hbar/\hbar$. The covariance before this change of units retains
$\hbar$; we express the admissibility criterion of Section~\ref{sec:covariance}
in dimensionless coordinates. With $\hbar=1$ and $t=s^2$, the original covariance is
$D(t,t^{-1})$, realizing the diagonal equality cases above.

\begin{corollary}\label{cor:physical}
Let $m,\omega,\hbar>0$ and set $s=\sqrt{\hbar/(m\omega)}$. Then
\[
 \gamma_{\psi_s}=\diag\left(\frac{\hbar}{m\omega},\hbar m\omega\right).
\]
\end{corollary}
\begin{proof}
Positivity gives $s>0$ and $s^2=\hbar/(m\omega)$. Substitution into
Theorem~\ref{thm:covariance} gives the formula.
\end{proof}
